\documentclass[11pt]{amsart}
\usepackage{amsmath,amssymb,amsthm}
\usepackage{graphicx}
\usepackage[colorlinks=true,linkcolor=blue,citecolor=blue,urlcolor=blue]{hyperref}

\newtheorem{theorem}{Theorem}
\newtheorem{lemma}{Lemma}

\newcommand{\R}{\mathbb R}
\newcommand{\N}{\mathbb N}
\newcommand{\cB}{\mathcal B}
\newcommand{\cK}{\mathcal K}
\newcommand{\wc}{\mathrm{wc}}
\newcommand{\dd}{\,d}

\title{Positive Local Compactness and Universal Hermite Projections}
\author{Agent lane 00}
\date{2026-07-03}

\begin{document}
\maketitle

\section*{Source and Target}

The source is Svela--Trapasso, \emph{A quantum harmonic analysis approach to nonlinear time-frequency concentration}, arXiv:2605.28786 \cite{ST}.  The authors formulate a broad open classification problem for smoothing and escape mechanisms of Cohen-class windows \(S\in \cB(L^2(\R^d))\).  The relevant introductory passage is shown below.

\begin{center}
\includegraphics[width=0.96\linewidth]{figures/open_problem_crop.png}
\end{center}

The source proves several structural criteria.  The first one used here is its positive localized-average criterion: if \(S\ge0\) and
\[
H_{\Omega,S}:=\chi_\Omega\star S=\int_\Omega \alpha_z(S)\dd z
\]
is compact, then \(f\mapsto Q_Sf\) is weak-to-norm continuous from \(L^2(\R^d)\) into \(L^p(\Omega)\), for every finite \(p\).  We first show that, for positive \(S\), this local compactness condition is also necessary.  The source also proves the endpoint formula \(\Lambda_{\infty,\Omega}(S)=w(S)\), where \(w(S)\) is the numerical radius, and numerical-radius attainment implies \(L^\infty\)-concentration attainment.  Near the end of the same discussion, the source notes that sparse Hermite functions give noncompact positive diagonal examples for a fixed bounded observation set.  The result below strengthens that observation to one Hermite projection in each dimension working simultaneously for every finite-measure observation window.

\section*{Proof Intuition}

For positive windows the source's sufficient local compactness mechanism is exact.  Indeed, if \(S\ge0\) and \(S\in\cB_{\wc}(p,\Omega)\), then every normalized weakly-null sequence \(f_n\) satisfies \(\|Q_Sf_n\|_{L^p(\Omega)}\to0\).  Since \(Q_Sf_n\ge0\), this forces
\[
\langle H_{\Omega,S}f_n,f_n\rangle=\int_\Omega Q_Sf_n\to0.
\]
A positive Hilbert-space operator with this vanishing property on all weakly-null sequences must be compact; otherwise an infinite-dimensional spectral subspace above a positive level supplies an orthonormal weakly-null countersequence.

The Hermite example then becomes an explicit noncompact positive window with compact localized average for every finite-measure \(\Omega\).  The source's fixed-window Hermite argument only needs local decay of the ambiguity functions on the difference set \(\Omega-\Omega\).  In dimension \(d\), take
\[
u_n^{(d)}=\phi_n\otimes\phi_0\otimes\cdots\otimes\phi_0 .
\]
The ambiguity function factors.  Its first factor tends to zero away from \(0\in\R^2\), so the product tends to zero except on the codimension-two set \(\{z_1=0\}\subset\R^{2d}\), which has measure zero.  Dominated convergence gives local \(L^2\)-decay on every ball.  Diagonalizing over all balls then makes \(H_{\Omega,S_d}\) compact for every bounded \(\Omega\).

If \(\Omega\) has finite measure but is unbounded, truncate it by \(\Omega_R=\Omega\cap B_R\).  The bounded pieces give compact operators \(H_{\Omega_R,S_d}\), while
\[
\|H_{\Omega,S_d}-H_{\Omega_R,S_d}\|
\le |\Omega\setminus B_R|\,\|S_d\|\to0 .
\]
Thus \(H_{\Omega,S_d}\) is a norm limit of compact operators.  The source's localized-average criterion then gives weak-to-norm continuity and hence finite-\(p\) attainment by the direct method.  At \(p=\infty\), the same projection attains because every nonzero orthogonal projection attains its numerical radius.

\section*{Definitions}

For \(f\in L^2(\R^d)\) and \(S\in\cB(L^2(\R^d))\), write
\[
Q_Sf(z)=\langle \alpha_z(S)f,f\rangle,\qquad z\in\R^{2d},
\]
for the Cohen-class distribution in the quantum harmonic analysis formulation of \cite{ST}.  For measurable \(0<|\Omega|<\infty\), the class \(\cB_{\wc}(p,\Omega)\) consists of those windows for which
\[
f_n\rightharpoonup f
\quad\Longrightarrow\quad
\|Q_Sf_n\|_{L^p(\Omega)}\to \|Q_Sf\|_{L^p(\Omega)}.
\]
The source proves that, for finite \(p\), this agrees with weak-to-norm continuity of \(f\mapsto Q_Sf\) into \(L^p(\Omega)\).

\section*{Theorem A: positive windows}

Let \(S\in\cB(L^2(\R^d))\) be positive, let \(\Omega\subset\R^{2d}\) be measurable with \(0<|\Omega|<\infty\), and let \(1\le p<\infty\).  Then
\[
S\in\cB_{\wc}(p,\Omega)
\quad\Longleftrightarrow\quad
H_{\Omega,S}:=\int_\Omega\alpha_z(S)\dd z\in\cK .
\]
Equivalently, for positive windows, weak-to-norm continuity over \(\Omega\) is exactly compactness of the mixed-state localized average.

\section*{Proof of Theorem A}

The implication \(H_{\Omega,S}\in\cK\Rightarrow S\in\cB_{\wc}(p,\Omega)\) is the source paper's positive localized-average criterion.

Conversely, assume \(S\in\cB_{\wc}(p,\Omega)\), and let \(f_n\rightharpoonup0\) be a normalized sequence in \(L^2(\R^d)\).  Since \(S\ge0\), also \(\alpha_z(S)\ge0\) and hence \(Q_Sf_n(z)\ge0\).  By weak-to-norm continuity,
\[
\|Q_Sf_n\|_{L^p(\Omega)}\to0.
\]
Because \(|\Omega|<\infty\), Hölder's inequality gives \(\|Q_Sf_n\|_{L^1(\Omega)}\to0\).  Therefore
\[
\langle H_{\Omega,S}f_n,f_n\rangle
=\int_\Omega \langle \alpha_z(S)f_n,f_n\rangle\dd z
=\int_\Omega Q_Sf_n(z)\dd z
\longrightarrow0
\]
for every normalized weakly-null sequence \(f_n\).

The operator \(H_{\Omega,S}\) is positive.  We claim that a positive bounded operator \(H\) on a Hilbert space is compact whenever \(\langle Hf_n,f_n\rangle\to0\) for every normalized weakly-null sequence.  If \(H\) were not compact, then by the spectral theorem there would be an \(\varepsilon>0\) such that the spectral projection \(1_{[\varepsilon,\infty)}(H)\) has infinite-dimensional range.  Choose an orthonormal sequence \(e_n\) in that range.  Then \(e_n\rightharpoonup0\), while
\[
\langle He_n,e_n\rangle\ge\varepsilon
\]
for all \(n\), a contradiction.  Thus \(H_{\Omega,S}\) is compact.

\section*{Theorem B: universal Hermite projections}

For every integer \(d\ge1\), there exists a noncompact positive orthogonal projection
\[
S_d=\sum_{k=1}^\infty u_k^{(d)}\otimes u_k^{(d)}
\quad\text{on }L^2(\R^d)
\]
with the following property: for every measurable \(\Omega\subset\R^{2d}\) with \(0<|\Omega|<\infty\) and every \(1\le p<\infty\),
\[
f_n\rightharpoonup f \text{ in } L^2(\R^d)
\quad\Longrightarrow\quad
\|Q_{S_d}f_n-Q_{S_d}f\|_{L^p(\Omega)}\to0 .
\]
Consequently \(S_d\in\cB_{\wc}(p,\Omega)\), and the concentration supremum
\[
\Lambda_{p,\Omega}(S_d)=\sup_{\|f\|_2=1}\|Q_{S_d}f\|_{L^p(\Omega)}
\]
is attained for every finite-measure observation set \(\Omega\).  In addition,
\[
\Lambda_{\infty,\Omega}(S_d)=\sup_{\|f\|_2=1}\|Q_{S_d}f\|_{L^\infty(\Omega)}
\]
is attained for every measurable \(\Omega\subset\R^{2d}\) with positive measure.  Thus these projections attain the concentration problem for every \(1\le p\le\infty\), while the weak-to-norm continuity conclusion is only claimed for finite \(p\).

\section*{Proof of Theorem B}

Let \((\phi_n)_{n\ge0}\) be the normalized one-dimensional Hermite basis.  For \(d\ge1\), define
\[
v_n^{(d)}=\phi_n\otimes\underbrace{\phi_0\otimes\cdots\otimes\phi_0}_{d-1\text{ copies}}
\in L^2(\R^d).
\]
These vectors are orthonormal as \(n\) varies.  Write \(z=(z_1,\ldots,z_d)\in(\R^2)^d=\R^{2d}\), where each \(z_j\in\R^2\).  The tensor-product structure of the Schrödinger representation gives
\[
|A(v_n^{(d)})(z)|
=|A(\phi_n)(z_1)|\prod_{j=2}^d |A(\phi_0)(z_j)|.
\]
In dimension one, the source records
\[
A(\phi_n)(\zeta)=e^{-\pi |\zeta|^2/2}L_n(\pi |\zeta|^2),
\]
and the consequences
\[
A(\phi_n)(\zeta)\to0\quad(\zeta\ne0),\qquad |A(\phi_n)(\zeta)|\le1.
\]
Therefore \(A(v_n^{(d)})(z)\to0\) for every \(z\) with \(z_1\ne0\).  The exceptional set \(\{z\in\R^{2d}:z_1=0\}\) has Lebesgue measure zero.  Since \(|A(v_n^{(d)})(z)|\le1\), dominated convergence yields, for every integer \(m\ge1\),
\[
a_{n,m}^{(d)}:=
\left(\int_{B_m^{2d}}|A(v_n^{(d)})(z)|^2\dd z\right)^{1/2}
\longrightarrow0
\qquad(n\to\infty),
\]
where \(B_m^{2d}\subset\R^{2d}\) is the centered radius-\(m\) ball.

Choose a strictly increasing subsequence \(n_k\) recursively so that
\[
a_{n_k,m}^{(d)}\le 2^{-k-m}\qquad\text{for every }1\le m\le k .
\]
This is possible because only finitely many conditions are imposed at stage \(k\).  Set \(u_k^{(d)}=v_{n_k}^{(d)}\), \(P_k=u_k^{(d)}\otimes u_k^{(d)}\), and define the strong operator sum
\[
S_d=\sum_{k=1}^\infty P_k .
\]
Since the \(u_k^{(d)}\)'s are orthonormal, \(S_d\) is the orthogonal projection onto their closed span.  It is positive, \(\|S_d\|=1\), and it is not compact because \(S_du_k^{(d)}=u_k^{(d)}\) while \(u_k^{(d)}\rightharpoonup0\).

First let \(E\subset\R^{2d}\) be bounded with \(0<|E|<\infty\).  Choose \(m\) such that \(E-E\subset B_m^{2d}\), and define
\[
C_E(u)^2=\int_E\int_E |A(u)(w-z)|^2\dd z\dd w .
\]
For \(k\ge m\), the diagonal choice gives
\[
\begin{aligned}
C_E(u_k^{(d)})^2
&=\int_{\R^{2d}}|A(u_k^{(d)})(v)|^2\,|E\cap(E-v)|\dd v\\
&\le |E|\int_{E-E}|A(u_k^{(d)})(v)|^2\dd v
\le |E|\int_{B_m^{2d}}|A(u_k^{(d)})(v)|^2\dd v\\
&\le |E|\,2^{-2k-2m}.
\end{aligned}
\]
Therefore \(\sum_k C_E(u_k^{(d)})<\infty\).

Let \(H_{E,k}=\chi_E\star P_k=\int_E\alpha_z(P_k)\dd z\).  The source computes
\[
\|H_{E,k}\|_{\mathcal S^2}=C_E(u_k^{(d)}),
\]
so \(\|H_{E,k}\|_{\cB}\le C_E(u_k^{(d)})\).  Since \(H_{E,k}\) is compact for every \(k\), the series
\[
H_{E,S_d}=\int_E\alpha_z(S_d)\dd z=\sum_{k=1}^\infty H_{E,k}
\]
converges in operator norm to a compact operator.  Thus \(H_{E,S_d}\in\cK\) for every bounded finite-measure \(E\).

Now let \(\Omega\subset\R^{2d}\) be arbitrary with \(0<|\Omega|<\infty\).  Put \(\Omega_R=\Omega\cap B_R^{2d}\).  Each \(H_{\Omega_R,S_d}\) is compact by the bounded case.  Moreover, for every measurable \(F\subset\R^{2d}\) of finite measure,
\[
\|H_{F,S_d}\|
=\left\|\int_F\alpha_z(S_d)\dd z\right\|
\le \int_F\|\alpha_z(S_d)\|\dd z
=|F|\,\|S_d\| .
\]
Consequently
\[
\|H_{\Omega,S_d}-H_{\Omega_R,S_d}\|
=\|H_{\Omega\setminus B_R^{2d},S_d}\|
\le |\Omega\setminus B_R^{2d}|\,\|S_d\|\longrightarrow0 .
\]
The operator \(H_{\Omega,S_d}\) is therefore a norm limit of compact operators, hence compact.

By Theorem A, or equivalently by the source's positive localized-average criterion applied to the compact operator \(H_{\Omega,S_d}\), this positive \(S_d\) satisfies, for every \(1\le p<\infty\),
\[
f_n\rightharpoonup f
\quad\Longrightarrow\quad
\|Q_{S_d}f_n-Q_{S_d}f\|_{L^p(\Omega)}\to0 .
\]
This proves \(S_d\in\cB_{\wc}(p,\Omega)\) for every finite-measure \(\Omega\).

For finite \(p\), weak-to-norm continuity gives attainment by the direct method.  The closed unit ball of \(L^2(\R^d)\) is weakly compact, \(f\mapsto\|Q_{S_d}f\|_{L^p(\Omega)}\) is weakly continuous on it, and \(S_d\ne0\) makes the maximum positive.  Since \(Q_{S_d}\) is homogeneous of degree two, every nonzero maximizer over the closed unit ball can be normalized to the unit sphere.  Thus \(\Lambda_{p,\Omega}(S_d)\) is attained for \(1\le p<\infty\).

It remains only to record the endpoint \(p=\infty\).  The source proves that \(\Lambda_{\infty,\Omega}(S)=w(S)\) for every positive-measure \(\Omega\), and that numerical-radius attainment of \(S\) implies endpoint concentration attainment.  Since \(S_d\) is an orthogonal projection, \(w(S_d)=\|S_d\|=1\), and the numerical radius is attained at every unit vector in \(\operatorname{Ran}S_d\).  Hence \(\Lambda_{\infty,\Omega}(S_d)=1\) is attained for every positive-measure \(\Omega\).

\section*{Verification and Novelty Check}

The proof is noncomputational.  It uses the source paper's positive localized-average criterion, its \(L^\infty\) numerical-radius formula, the Hermite ambiguity decay quoted in the source, and the elementary tensor factorization of ambiguity functions.  The new structural step is the converse in Theorem A: for positive windows, compactness of \(H_{\Omega,S}\) is necessary as well as sufficient for finite-\(p\) weak continuity.  The new Hermite steps are the all-dimensional tensor-Hermite local decay observation, the simultaneous diagonal choice over all balls, and the finite-measure truncation argument for \(H_{\Omega,S_d}\).

Cheap run indexes were searched for arXiv:2605.28786, the exact title, `positive localized average', `compact localized average', `Hermite projection', `positive diagonal', `finite measure', `all dimensions', and related phrases.  No existing packet or attempt for this strengthened result was found.  A bounded web search for the exact source title, positive localized-average converse phrases, and strengthened Hermite finite-measure phrases found the arXiv source record and no later answer.

\section*{Limitations}

This is not a full classification of all Cohen-class windows.  It does give an exact finite-\(p\) weak-continuity classification for positive bounded windows, and one explicit noncompact positive projection in each dimension that works for every finite-measure observation window in that dimension.  It does not classify arbitrary nonpositive windows.  The endpoint \(p=\infty\) is covered for attainment of the Hermite projections only, not by weak-to-norm continuity.

\section*{Human Review Recommendation}

First review Theorem A, especially the spectral-theorem argument that positivity plus vanishing quadratic forms on normalized weakly-null sequences forces compactness.  Then review the tensor factorization of the ambiguity functions and the dominated-convergence step on balls in \(\R^{2d}\).  Then review the estimate of \(C_E(u_k^{(d)})\) for bounded \(E\), the norm convergence of \(H_{E,S_d}=\sum_kH_{E,k}\), and the truncation step
\[
H_{\Omega,S_d}=\lim_{R\to\infty}H_{\Omega\cap B_R^{2d},S_d}
\]
for unbounded finite-measure \(\Omega\).  If these are accepted, the finite-\(p\) result follows directly from the source paper's positive localized-average criterion.  For the endpoint, check that \(S_d\) attains its numerical radius as an orthogonal projection, and apply the source paper's \(L^\infty\) numerical-radius formula.

\begin{thebibliography}{9}
\bibitem{ST}
E. A. T. Svela and S. I. Trapasso,
\emph{A quantum harmonic analysis approach to nonlinear time-frequency concentration},
arXiv:2605.28786, 2026.
\end{thebibliography}

\end{document}
